\section{Conclusion and Future Scope}
We presented INN-RC, a lossless steganography scheme in which an integer coupling network is conditioned on the
message instead of taking it as an input. Because every coupling step is an integer addition driven by the other pixel
half and the message, the cover is recovered exactly once the message is known, and the only residual that has to be
compensated is the set of wrongly decoded bits and the clipped pixels. Coding this residual with contexts the receiver
can rebuild (decoder confidence and the re-estimated coupling offset) keeps it smaller than the message itself.

\textbf{Exact reversibility.} In all \exactQ{} successful IC-INN-Q runs and \exactR{} IC-INN-R runs on \nImgs{} test
images, the recovered cover was identical to the original pixel for pixel and the extracted message had BER $=0$
(NC $=1$), which we verified with \texttt{isequal} in MATLAB rather than with a PSNR threshold.

\textbf{Limitations.} (i) On a guaranteed lossless channel classical PEE still gives a higher PSNR at these payloads;
our scheme pays for robustness with image quality. (ii) At 4096 bits a few busy or noisy dark images cannot hold the
side information and embedding is refused. (iii) The receiver must reproduce the rounded coupling offsets exactly, so
sender and receiver should run the same floating-point code; a fixed-point network would remove this assumption.
(iv) We work on 8-bit grayscale images and converted the DICOM scans to 8 bits, and BOSSbase could not be included.

\textbf{Future work.} Fixed-point (integer-only) coupling networks for platform-independent decoding, native 12/16-bit
DICOM and colour images, a learned predictor for the PEE stage to increase its capacity, larger payloads through a
multi-scale message map, and testing against deep steganalysers.
